\section{Introduction}

Electroencephalography (EEG) datasets vary simultaneously in montage, channel
count, sampling rate, recording length, subject population, and label space.
EEG foundation models address this heterogeneity with large-scale
self-supervised pretraining and specialized tokenization or attention
mechanisms \citep{yang2023biot,jiang2024labram,wang2025cbramod,elouahidi2025reve}, often borrowing
masked-reconstruction and contrastive objectives from vision. Their success
has made ``more EEG pretraining'' a natural explanation for improved transfer,
yet the downstream architecture and the input geometry are usually changed at
the same time, making it hard to separate domain pretraining from a better
representation interface.

This motivates our central question: \emph{is vision pretraining a free lunch
for EEG decoding?}  We use ``free lunch'' in a restricted sense: reusing
publicly available visual weights without collecting EEG pretraining data or
running an EEG self-supervised stage. It does not mean zero computation,
universal improvement, or deployment without task-specific fine-tuning.

We study this confound by comparing EEG-pretrained foundation models with
vision-pretrained convolutional and Transformer backbones that use no
EEG-specific pretraining. A raw EEG segment is naturally a matrix with
electrodes as rows and time samples as columns. Passing it directly to a 2D
CNN is unattractive: the input is extremely wide, often only 6--64 rows high,
and standard vision networks reduce both spatial axes by a factor near 32, so
the electrode dimension can collapse after a few stages while most spatial
computation traverses a long time axis. Resizing would change values, and a
spectrogram would discard phase and introduce analysis choices. The problem is
most severe for sparse montages (4--16 channels), which need a geometry
adapter to preserve signal-bearing height and expose structured temporal
patterns to local 2D operators; dense 64-channel montages can remain unfolded.

Our alternative is a lossless index transformation.  For fold factor $P$, we
reshape each channel into groups of $P$ consecutive samples and transpose the
within-group phase into the row axis, mapping $C\times T$ to
$(CP)\times(T/P)$.  We call this operation \emph{temporal folding}.  It has no
learned parameters and is exactly invertible.  However, it changes what a
local convolution means: a vertical step within an electrode advances by one
sample, whereas a horizontal step advances by $P$ samples.  A small 2D kernel
therefore combines fine temporal differences with $P$-spaced context, while
the taller representation carries more nonzero rows through the height
hierarchy.

We evaluate this interface across heterogeneous EEG tasks and analyze when it
helps or fails, using EfficientNet-B0, ConvNeXt-Tiny, and ViT-Small
across 12 downstream benchmarks, keeping conventional supervised models and
the BIOT, LaBraM, CBraMod, and REVE foundation models as contextual baselines.

Our contributions are:

\begin{itemize}
    \item \textbf{A lossless EEG--vision interface.} Temporal folding is a
    bijective channel--time reindexing with no interpolation, filtering, or
    learned parameters, preserving signal-bearing height for sparse montages.
    \item \textbf{Evidence that visual representations transfer to EEG.}
    ImageNet- and DINOv3-initialized CNN and Transformer backbones learn from
    the same folded waveform interface across heterogeneous tasks.
    \item \textbf{A broad benchmark comparison.} We compare the visual
    transfer with supervised architectures and EEG foundation models, including
    CBraMod and REVE, on 12 datasets under shared benchmark splits.
\end{itemize}

Our central claim is specific and testable: when EEG is mapped to a geometry
that a vision encoder can process without discarding temporal information,
mature visual representations become a competitive cross-domain prior for EEG
decoding. The effect is substantial but not universal—EEG-specific foundation
models remain ahead on several tasks—so this is not a replacement claim. It is
a design principle: before investing in larger EEG-specific pretraining,
quantify the gain obtained from geometry alignment and established
representations from other modalities.
